\documentclass[11pt]{article}
\usepackage[margin=1in]{geometry}
\usepackage{amsmath}
\usepackage{booktabs}
\usepackage{hyperref}

\title{RADCAL Neural-Network Database Index Mapping}
\author{}
\date{}

\begin{document}
\maketitle

\section{How to Build the FORTRAN CODE}

The neural-network database generator is built from
\texttt{rmod.f90}, \texttt{rcal.f90}, and \texttt{nn\_database.f90}. The
resulting executable is named \texttt{radcal\_nn}.

\subsection{macOS}

From the repository root, run:

\begin{verbatim}
Build/intel_osx_64/make_radcal_nn.sh
\end{verbatim}

The script selects the available macOS Fortran compiler and invokes the
\texttt{intel\_osx\_64\_nn} Makefile target. The executable is written to:

\begin{verbatim}
Build/intel_osx_64/radcal_nn
\end{verbatim}

\subsection{Linux}

From the repository root, run:

\begin{verbatim}
Build/intel_linux_64/make_radcal_nn.sh
\end{verbatim}

The script configures the Intel Fortran compiler and invokes the
\texttt{intel\_linux\_64\_nn} Makefile target. The executable is written to:

\begin{verbatim}
Build/intel_linux_64/radcal_nn
\end{verbatim}

\subsection{Cleanup}

The clean target removes \texttt{*.o}, \texttt{*.mod}, and \texttt{*.obj}.
It does not remove the \texttt{radcal\_nn} executable.

\section{Generate NN Database}

The database is generated by
\path{Examples/NN_TEST/RADCAL_TRAINING_DATA_GENERATION/generate_nn_database.py}.
The script launches the
\texttt{radcal\_nn} Fortran executable, supplies the namelist filename and
output prefix, and reports progress as each temperature is completed.

First change to the training-data generation directory:

\begin{verbatim}
cd Examples/NN_TEST/RADCAL_TRAINING_DATA_GENERATION
\end{verbatim}

Then generate the test database on macOS with:

\begin{verbatim}
python3 generate_nn_database.py \
  ../../../Build/intel_osx_64/radcal_nn \
  --input RADCAL_NN_TEST.IN \
  --output-prefix radcal_nn_test \
  --processes 4
\end{verbatim}

On Linux, replace \texttt{intel\_osx\_64} with
\texttt{intel\_linux\_64}.

Each MPI rank evaluates a different subset of temperatures and writes its
results directly to the correct offsets in the shared Kappa file. A single
process uses the same MPI code path. If \texttt{--processes} exceeds the
\texttt{N\_RTMP} value in the input namelist, the Python launcher
automatically reduces it to \texttt{N\_RTMP}. The
header, valid-gas index file, Kappa file layout, and Python reader are
unchanged by the number of processes.
When the requested process count exceeds the CPU slots available to Open MPI,
add \texttt{--oversubscribe} to the Python command.

\subsection{Input namelist}

The input file
\path{Examples/NN_TEST/RADCAL_TRAINING_DATA_GENERATION/RADCAL_NN_TEST.IN}
uses the
\texttt{NN\_INPUT} namelist. It specifies the minimum and maximum gas mole
fractions, minimum and maximum soot volume fractions, number of points on each
axis, temperature limits, number of temperature points, and optionally the
pressure. Pressure defaults to 1~atm when it is omitted.

The gas and soot axes contain zero at index 0, followed by logarithmically
spaced positive values from the configured minimum through the maximum.
Temperature is linearly spaced. Nitrogen is the background gas:
\[
x_{\mathrm{N_2}} =
1-x_{\mathrm{CO_2}}-x_{\mathrm{H_2O}}-x_{\mathrm{CO}}-x_{\mathrm{C_2H_4}}.
\]
Gas combinations whose participating-species mole fractions sum to more than
one are discarded before RADCAL is called.

\subsection{Output files}

The output prefix produces three files under
\path{Examples/NN_TEST/RADCAL_TRAINING_DATA_GENERATION}:

\begin{center}
\begin{tabular}{ll}
\toprule
File & Contents \\
\midrule
\texttt{radcal\_nn\_test\_header.txt} & Configuration, dimensions, and axes \\
\texttt{radcal\_nn\_test\_valid\_gas\_indices.bin} & Valid gas-index combinations \\
\texttt{radcal\_nn\_test\_kappa.bin} & Planck and 10~cm Kappa values \\
\bottomrule
\end{tabular}
\end{center}

The readable header records the physical configuration, binary data types,
record order, index equation, and lookup tables for temperature, gas mole
fraction, and soot volume fraction. Binary values use the native byte order
listed in the header.

\subsection{How indices work in the output files}

All indices are zero-based. The gas-axis table in the header maps a gas-axis
index \(i\) to mole fraction \(x_i\). Each record in
\texttt{radcal\_nn\_test\_valid\_gas\_indices.bin} contains four native
32-bit integers:
\[
(i_{\mathrm{CO_2}},i_{\mathrm{H_2O}},i_{\mathrm{CO}},i_{\mathrm{C_2H_4}}).
\]
This record represents
\[
(x_{i_{\mathrm{CO_2}}},x_{i_{\mathrm{H_2O}}},
x_{i_{\mathrm{CO}}},x_{i_{\mathrm{C_2H_4}}}).
\]
The position of the record in the file is the valid-gas index \(g\). Each
record occupies 16 bytes, so its byte offset is
\[
\operatorname{gas\_byte\_offset}(g)=16g.
\]

Each record in \texttt{radcal\_nn\_test\_kappa.bin} contains two native
32-bit floating-point values:
\[
(\kappa_{\mathrm{Planck}},\kappa_{10\,\mathrm{cm}}).
\]
Soot index changes fastest, valid-gas index changes next, and temperature
index changes slowest. Let \(t\) be the temperature index, \(g\) the
valid-gas index, \(s\) the soot index, \(N_g\) the number of valid gas
compositions, and \(N_s\) the number of soot values. The Kappa record index is
\[
r(t,g,s)=\left(tN_g+g\right)N_s+s,
\]
and its byte offset is
\[
\operatorname{kappa\_byte\_offset}(t,g,s)=8r(t,g,s).
\]

For the retained reduced MPI fixture generated from
\texttt{RADCAL\_NN\_PROGRESS\_TEST.IN}, \(N_g=20\) and \(N_s=3\). For
\(t=1\), \(g=3\), and \(s=1\),
\[
r=(1\cdot20+3)\cdot3+1=70,
\]
so the Kappa record begins at byte \(8\cdot70=560\).

\section{Read NN Database}

The memory-mapped reader is
\path{Examples/NN_TEST/VERIFICATION_TRAINING_DATA_WITH_FDS/read_nn_database.py}.
It reads the header,
memory-maps both binary files, reconstructs physical input values from their
indices, and returns batches without loading the complete database into
memory.

The constructor \texttt{\_\_init\_\_} runs when a
\texttt{RadcalNNDataset} object is created. It reads the dimensions and axis
tables from the header, resolves the binary-file paths, and opens the gas-index
and Kappa files with \texttt{numpy.memmap}. A memory map behaves like a NumPy
array, but the operating system reads only the portions of the files that are
accessed.

\subsection{Record reconstruction}

The variable \texttt{size} is the total number of database records:

\[
N_{\mathrm{records}}=N_T\,N_g\,N_s,
\]

where \(N_T\), \(N_g\), and \(N_s\) are the numbers of temperatures, valid gas
compositions, and soot values. For a flattened record number \(r\), the reader
recovers the three database indices using

\[
s=r\bmod N_s,\qquad
g=\left\lfloor\frac{r}{N_s}\right\rfloor\bmod N_g,\qquad
t=\left\lfloor\frac{r}{N_sN_g}\right\rfloor.
\]

Thus soot changes fastest, gas composition changes next, and temperature
changes slowest. The gas-index record then maps \(g\) to the four gas-axis
indices for \(\mathrm{CO_2}\), \(\mathrm{H_2O}\), \(\mathrm{CO}\), and
\(\mathrm{C_2H_4}\).

\subsection{Inputs, targets, and batches}

The method \texttt{batch(start,count)} returns two two-dimensional arrays.
For a batch containing \(B\) records, \texttt{inputs} has shape \(B\times7\),
with columns temperature, pressure, \(\mathrm{CO_2}\), \(\mathrm{H_2O}\),
\(\mathrm{CO}\), soot volume fraction, and \(\mathrm{C_2H_4}\).
\texttt{targets} has shape \(B\times2\), with Planck-mean Kappa in column 0
and 10~cm Kappa in column 1.

The method \texttt{batches(batch\_size)} uses Python's \texttt{yield} statement
to return one batch at a time. This allows neural-network training to loop over
a database much larger than the available memory.

Inspect a batch from the command line with:

\begin{verbatim}
python3 Examples/NN_TEST/VERIFICATION_TRAINING_DATA_WITH_FDS/\
read_nn_database.py \
  Examples/NN_TEST/RADCAL_TRAINING_DATA_GENERATION/\
radcal_nn_test_header.txt \
  --start 0 --count 8
\end{verbatim}

Use the reader in Python with:

\begin{verbatim}
from Examples.NN_TEST.VERIFICATION_TRAINING_DATA_WITH_FDS.\
read_nn_database import RadcalNNDataset

data = RadcalNNDataset(
    "Examples/NN_TEST/RADCAL_TRAINING_DATA_GENERATION/"
    "radcal_nn_test_header.txt"
)

for inputs, targets in data.batches(batch_size=65536):
    model.train_on_batch(inputs, targets)
\end{verbatim}

The input columns are temperature, pressure, \(\mathrm{CO_2}\),
\(\mathrm{H_2O}\), \(\mathrm{CO}\), soot volume fraction, and
\(\mathrm{C_2H_4}\). The target columns are Planck mean Kappa and 10~cm
path-length-specific Kappa.

\section{Direct Database Verification with RADCAL}

The direct verification program is
\path{Examples/NN_TEST/VERIFICATION_TRAINING_DATA_WITH_FDS/%
generate_direct_verification.f90}. It calls RADCAL directly using the same
10~cm path length, pressure, wall temperature, spectral limits, and
discretization rules as the database generator. The program reads
\texttt{RADCAL\_NN\_TEST.IN} and writes small text files for the following
checks:

\begin{enumerate}
   \item Planck-mean Kappa versus the complete temperature axis for pure
         \(\mathrm{CO_2}\), \(\mathrm{H_2O}\), \(\mathrm{CO}\), and
         \(\mathrm{C_2H_4}\), and for soot near \(f_v=10^{-5}\).
   \item 10~cm Kappa versus the complete gas or soot axis at the temperature
         nearest 1527~K, with all other species set to zero.
   \item At the temperature nearest 1527~K, 10~cm Kappa versus
         \(\mathrm{H_2O}\) with
         \(x_{\mathrm{CO_2}}\) fixed near
         \(8.1854673070690193\times10^{-2}\).
   \item At the temperature nearest 1527~K, 10~cm Kappa versus soot volume
         fraction with \(x_{\mathrm{CO_2}}\) fixed near
         \(8.1854673070690193\times10^{-2}\) and
         \(x_{\mathrm{H_2O}}\) fixed near
         \(4.9619476030029037\times10^{-2}\).
\end{enumerate}

The requested temperature and fixed mole fractions are mapped to the nearest
values on the configured database axes. For the 22-temperature and 25-point
configuration, these values are approximately 1527.14286~K,
\(x_{\mathrm{CO_2}}=8.18546731\times10^{-2}\), and
\(x_{\mathrm{H_2O}}=4.96194760\times10^{-2}\). The direct results are converted
to the same 32-bit floating-point representation stored in the binary Kappa
file.

\subsection{Build and generate direct results}

From the repository root, run:

\begin{verbatim}
cd Examples/NN_TEST/VERIFICATION_TRAINING_DATA_WITH_FDS
chmod +x build_direct_verification.sh
./build_direct_verification.sh
./generate_direct_verification \
  ../RADCAL_TRAINING_DATA_GENERATION/RADCAL_NN_TEST.IN .
\end{verbatim}

The default build uses \texttt{gfortran}. To use another compiler, pass its
command to the build script. For example:

\begin{verbatim}
./build_direct_verification.sh ifort
\end{verbatim}

Using the same compiler options as the database build minimizes numerical
differences before conversion to 32-bit storage.

\subsection{Plot and compare with the database}

Run the Python comparison from the same directory:

\begin{verbatim}
python3 compare_direct_verification.py \
  ../RADCAL_TRAINING_DATA_GENERATION/radcal_nn_test_header.txt
\end{verbatim}

The reader locates the gas-index and Kappa binary files relative to the header.
For every direct result, it reconstructs the matching temperature, valid-gas,
and soot indices and reads the corresponding database value.

The comparison script first creates all 12 overlay plots in
\texttt{direct\_comparison\_plots}. Direct RADCAL values are shown as solid
lines, and database values are shown as open-circle markers. Only after every
plot has been written does the script report the complete PASS/FAIL summary.
Values are considered equal using a relative tolerance of \(10^{-6}\) and an
absolute tolerance of \(10^{-8}\):
\[
|a-b|\leq 10^{-8}+10^{-6}|b|.
\]
If one or more comparisons exceed this tolerance, a final assertion lists
every failed case and its maximum absolute error instead of stopping at the
first mismatch.

To select a different plot directory, use:

\begin{verbatim}
python3 compare_direct_verification.py \
  ../RADCAL_TRAINING_DATA_GENERATION/radcal_nn_test_header.txt \
  --plot-directory my_verification_plots
\end{verbatim}

\subsubsection{FDS-style K10 lookup comparison}

The program \texttt{generate\_pl10cm\_lookup.f90} compares a lookup-table
estimate of the 10~cm absorption coefficient, \(\kappa_{10\,\mathrm{cm}}\),
with the direct RADCAL values already stored by \texttt{nn\_database.f90}.
The direct values are not recomputed by this comparison. The program reads
each binary record as the Planck-mean coefficient followed by the direct
10~cm coefficient, and uses the latter as the reference value.

The lookup construction and query reproduce the relevant behavior in
\texttt{fds/Source/radi.f90}: 50 logarithmically spaced composition entries
stored in slots 1--50, a zero-valued slot 0 for below-range composition,
45 temperature slices selected by floor-and-clamp, linear interpolation in
composition only. For a like-for-like comparison with \texttt{nn\_database.f90},
the table stores \texttt{AMEAN} directly rather than FDS's
\texttt{MIN(AMEAN,AP0)} value. The comparison retains \(\mathrm{cm^{-1}}\)
units because the NN database is stored in these units; FDS converts its
completed lookup table to \(\mathrm{m^{-1}}\) for use in the solver.

Build and run the comparison from the verification directory:

\begin{verbatim}
chmod +x build_pl10cm_lookup.sh
./build_pl10cm_lookup.sh
./generate_pl10cm_lookup \
  ../RADCAL_TRAINING_DATA_GENERATION/radcal_nn_test_header.txt \
  ../RADCAL_TRAINING_DATA_GENERATION/radcal_nn_test_valid_gas_indices.bin \
  ../RADCAL_TRAINING_DATA_GENERATION/radcal_nn_test_kappa.bin
\end{verbatim}

This writes \texttt{fds\_lookup\_comparison.csv}, which contains mean and
maximum absolute percent error at each temperature, and a bounded sample of
direct-versus-lookup pairs in \texttt{fds\_lookup\_scatter.csv}, to the
directory from which the program is run.

\paragraph{Temperature-interpolated option}

The program also writes
\texttt{fds\_lookup\_scatter\_temperature\_interpolated.csv}. It evaluates
the same composition interpolation as the FDS-style lookup, but linearly
interpolates between the two surrounding temperature slices. In contrast, the
FDS-style result uses the lower, floor-and-clamp temperature index. The
interpolated result is therefore an accuracy comparison option, not an
emulation of \texttt{GET\_KAPPA} in FDS. Comparing the two isolates the error
introduced by the 50~K table-temperature spacing.

Create the two linear scatter plots with:

\begin{verbatim}
python3 plot_fds_lookup_comparison.py fds_lookup_scatter.csv
\end{verbatim}

The resulting \texttt{fds\_lookup\_vs\_direct\_fds\_linear.png} and
\texttt{fds\_lookup\_vs\_direct\_temperature\_interpolated\_linear.png}
show the FDS floor-and-clamp lookup and the temperature-interpolated lookup,
respectively. The script automatically reads the interpolated scatter file
from the same directory. The diagonal indicates exact agreement between the
lookup result and the direct database value.

\end{document}
